%% LyX 2.4.4 created this file.  For more info, see https://www.lyx.org/.
%% Do not edit unless you really know what you are doing.
\documentclass[english]{article}
\usepackage[T1]{fontenc}
\usepackage[utf8]{inputenc}
\setcounter{secnumdepth}{-2}
\setlength{\parindent}{0bp}
\usepackage{amsmath}
\usepackage{amsthm}
\usepackage{amssymb}

\makeatletter
%%%%%%%%%%%%%%%%%%%%%%%%%%%%%% Textclass specific LaTeX commands.
\theoremstyle{plain}
\newtheorem*{thm*}{\protect\theoremname}

\@ifundefined{date}{}{\date{}}
%%%%%%%%%%%%%%%%%%%%%%%%%%%%%% User specified LaTeX commands.
\usepackage{graphicx}
\usepackage{geometry}
\newcommand{\limi}{\underset{n\rightarrow\infty}{\lim}}
\newcommand{\limxz}{\underset{x\rightarrow x_{0}}{\lim}}
\newcommand{\limfxz}{\underset{x\rightarrow x_{0}}{\lim}f(x)}
\newcommand{\limfxm}{\underset{x\rightarrow x_{0}^{-}}{\lim}f(x)}
\newcommand{\limfxp}{\underset{x\rightarrow x_{0}^{+}}{\lim}f(x)}
\newcommand{\limxm}{\underset{x\rightarrow x_{0}^{-}}{\lim}}
\newcommand{\limxp}{\underset{x\rightarrow x_{0}^{+}}{\lim}}
\newcommand{\sqrm}{M_{m\times m}(\mathbb{F})}
\newcommand{\seqa}{(a_{n})_{n=1}^{\infty}}
\newcommand{\seqx}{(x_{n})_{n=1}^{\infty}}
\newcommand{\funcdr}{f:D\rightarrow\mathbb{R}}
\newcommand{\der}{\limxz\frac{f(x)-f(x_{0})}{x-x_{0}}}
\usepackage[all]{pdfprivacy}

\makeatother

\usepackage{babel}
\providecommand{\theoremname}{Theorem}

\begin{document}
\title{{\Large Weierstrass Extreme Value Theorem\vspace{-2em}}}
\maketitle
\begin{thm*}
Let $f:\left[a,b\right]\rightarrow\mathbb{R}$ be a continuous function
from $\left[a,b\right]$ to $\mathbb{R}$.

Then $f$ attains a global maximum and a global minimum on $\left[a,b\right]$.
\[
\exists c,d\in\left[a,b\right]\ :\ \forall x\in\left[a,b\right]\quad f(c)\le f(x)\le f(d)
\]
\end{thm*}
\medskip{}

\rule[0.5ex]{1\columnwidth}{1pt}

\medskip{}

\begin{proof}
We will first prove that $f$ attains a global maximum.

According to the Weierstrass Boundedness Theorem, $f$ is bounded
on $\left[a,b\right]$.

Hence $\operatorname{Im}(f)$ is a bounded subset of $\mathbb{R}$
and it is also a non-empty set (since $f(a)\in\operatorname{Im}(f)$).

The upper bound theorem states that every non-empty subset of $\mathbb{R}$
which is bounded from above has a supremum, so $\operatorname{Im}(f)$
has a supremum. 

Denote $\beta=\sup\left(\operatorname{Im}(f)\right)$.

Also recall that since $\beta$ is a supremum, it fulfills the supremum
characterization 
\[
\forall\varepsilon>0\ \exists y\in\operatorname{Im}(f)\ :\ \beta-\varepsilon<y\le\beta
\]
And in particular,
\[
\forall n\in\mathbb{N}\ \exists y_{n}\in\operatorname{Im}(f)\ :\ \beta-\frac{1}{n}<y_{n}\le\beta
\]
\[
\therefore\ \forall n\in\mathbb{N}\ \exists x_{n}\in\left[a,b\right]\ :\ \beta-\frac{1}{n}<f(x_{n})\le\beta
\]
We shall examine $\left(x_{n}\right)_{n=1}^{\infty}$. 

It is a bounded sequence, therefore according to the Bolzano-Weierstrass
theorem, it has a convergent subsequence $\left(x_{n_{k}}\right)_{k=1}^{\infty}$.

Denote $\underset{k\rightarrow\infty}{\lim}x_{n_{k}}=d$. 

Since limits of sequences weakly preserve inequalities maintained
by all elements of the sequence, $d\in\left[a,b\right]$.

Because $f$ is continuous, Heine's characterization of continuous
functions dictates $\underset{k\rightarrow\infty}{\lim}f\left(x_{n_{k}}\right)=f(d)$.

Moreover, let us notice that since $n_{k}\ge k$ we get
\[
\forall k\in\mathbb{N},\ \beta\ge f\left(x_{n_{k}}\right)>\beta-\frac{1}{n_{k}}\ge\beta-\frac{1}{k}
\]

Since $\underset{k\rightarrow\infty}{\lim}\left(\beta-\frac{1}{k}\right)=\beta$,
the squeeze theorem implies 
\[
f(d)=\underset{k\rightarrow\infty}{\lim}f\left(x_{n_{k}}\right)=\beta
\]

Recall that a supremum which belongs to the set is always the maximum
of the set, hence $\beta=\max\left(\operatorname{Im}(f)\right)$ and
$d$ is a global maximum point in $\left[a,b\right]$, as needed.

As for the global minimum: we know that a scalar multiple of a continuous
function is also continuous, hence $-f$ is a continuous function,
and the global maximum result applies to it:
\[
\exists c\in\left[a,b\right]\ :\ \forall x\in\left[a,b\right]\quad-f(x)\le-f(c)
\]

\[
\therefore\quad\exists c\in\left[a,b\right]\ :\ \forall x\in\left[a,b\right]\quad f(x)\ge f(c)
\]

This gives us the global minimum result.
\end{proof}

\end{document}
